\documentclass[11pt]{article}
\usepackage[margin=1in]{geometry}
\usepackage{amsmath,amssymb}
\usepackage{graphicx}
\usepackage{booktabs}
\usepackage[colorlinks=true,linkcolor=blue,citecolor=blue,urlcolor=blue]{hyperref}

\title{Inhomogeneous-media benchmarks for \texttt{lebedev\_em}}
\author{Shari Moskow\thanks{Department of Mathematics, Drexel University.
Prepared with assistance from Claude (Anthropic).}}
\date{July 17, 2026}

\newcommand{\sdot}{\dot\sigma}

\begin{document}
\maketitle

\begin{abstract}
We document two inhomogeneous-media benchmarks for the \texttt{lebedev\_em}
solver, an implementation of the Lebedev staggered-grid finite-difference
scheme of Davydycheva, Druskin \& Habashy \cite{DDH03} (DDH03). The first is a
two-half-space model with a $10\times$ conductivity contrast, validated
against the exact Sommerfeld-integral solution: the four-cluster Lebedev
average attains $\approx 1\%$ accuracy in the conductive half-space, and the
cluster averaging reduces the single-cluster error by a factor of five. The
second reproduces the self-consistency convergence test of DDH03 Figure~7 ---
a borehole with invasion crossing a $60^\circ$ dipping anisotropic layer ---
where solutions on optimal geometric grids with $k=3,\dots,6$ transverse steps
are graphically indistinguishable, and the $k=5$ solution agrees with the
$k=6$ reference to within $1.5\%$ at every receiver. Both benchmarks were run
on the repository state that includes the July 2026 correctness fixes.
\end{abstract}

\section{Background}

The solver discretizes the frequency-domain Maxwell system (time convention
$e^{-i\omega t}$)
\begin{equation}
\nabla\times \mathbf{E} = i\omega\mu\mathbf{H}, \qquad
\nabla\times \mathbf{H} = \sdot\,\mathbf{E} + \mathbf{J},
\qquad \sdot = \sigma - i\omega\varepsilon,
\end{equation}
on Lebedev's staggered grid, which splits into four Yee-like clusters carrying
mixed electric/magnetic boundary conditions; averaging the four cluster
solutions cancels both the source-discretization and domain-truncation errors
(DDH03, \S``Superconvergence''). Heterogeneous and anisotropic media are
handled by sub-cell homogenization (sequential nodal homogenization, after
Moskow et al.~\cite{MDHLD99}), and the transverse directions use the optimal
geometric grids of Ingerman--Druskin--Knizhnerman with progression factor
$\alpha = \exp(\gamma\pi/\sqrt{k})$, $\gamma = 1/\sqrt2$.

Both benchmarks below exercise the full inhomogeneous pipeline: media
builders, homogenization at material interfaces, the four-cluster solve with
per-cluster mixed boundary conditions, and the interpolate-then-average field
extraction of DDH03 eq.~(7).

\section{Benchmark 1: two half-spaces, $10\times$ contrast}

\paragraph{Model.} A unit vertical electric dipole (VED) at the origin in
\[
\sigma(z) = \begin{cases} 0.1 \text{ S/m}, & z < 4\text{ m},\\
1.0 \text{ S/m}, & z \ge 4 \text{ m},\end{cases}
\qquad f = 2500 \text{ Hz},
\]
with receivers on the $z$-axis, $z\in[0.5,\,7.5]$ m, straddling the contact.
The reference is the exact Sommerfeld-integral solution for $E_z$ in a
two-layer conductor (implemented in \texttt{analytics.py} and independently
verified: it reduces to the closed-form homogeneous solution when
$\sigma_1=\sigma_2$ to $3\times10^{-14}$ relative, and satisfies the
normal-current continuity $\sigma_1 E_z^- = \sigma_2 E_z^+$ at the contact).

\paragraph{Grid.} Optimal geometric transverse grids with $k=3$
($M_x=M_y=12$), a hybrid axial grid with fine equidistant spacing
$\Delta z = 0.0625$ m near the source ($M_z=160$, $N_R = 13{,}604$,
$\approx 41$k complex degrees of freedom), piecewise-constant $\sigma$ with
the contact resolved by the layered media builder.

\paragraph{Results.} Figure~\ref{fig:twolayer} shows $|E_z|$ on the axis, the
pointwise deviation from the analytic solution, and the contact effect
(ratio of the two-layer to the homogeneous response).

\begin{itemize}
\item In the conductive half-space ($z \ge 4$ m), across the $10\times$ jump,
the RMS relative error of $\mathrm{Re}\,E_z$ is $\mathbf{1.0\%}$.
\item The overall RMS error ($5.6\%$) is dominated by the near-source nodes
$z \lesssim 1.7$ m, where the transverse domain is only $0.3\,\delta_1$ wide
($\delta_1 \approx 32$ m in the resistive layer); the \emph{same} error is
present in the homogeneous control run (RMS $5.7\%$), so the material
interface itself contributes essentially nothing --- the FD scheme crosses
the contact cleanly.
\item The sharp physical structure at the contact is captured quantitatively:
the two-layer/homogeneous ratio rises to $1.65$ just above the contact
(charge accumulation on the resistive side) and drops to $\approx 0.18$
below it, matching the analytic curve point by point
(Fig.~\ref{fig:twolayer}, right panel).
\item Cluster averaging matters: reading a single cluster (C000) gives RMS
$29.9\%$; the four-cluster Lebedev average reduces this to $5.7\%$ --- the
error-cancellation mechanism of DDH03 at work in an inhomogeneous medium.
\end{itemize}

\begin{figure}[htbp]
\centering
\includegraphics[width=\textwidth]{../examples/benchmark_two_layer.png}
\caption{Two-half-space benchmark ($\sigma_1=0.1$, $\sigma_2=1.0$ S/m,
contact at $z_c=4$ m, $f=2500$ Hz). Left: $|E_z|$ on the axis, FD (symbols)
vs.\ analytic (lines), homogeneous control in blue, two-layer in red.
Center: pointwise percent deviation from the analytic solutions. Right:
contact effect --- the ratio of two-layer to homogeneous $E_z$, FD (black)
vs.\ Sommerfeld analytic (red). The isolated spike at $z=4$~m in the center
panel sits at the contact node itself, where $E_z$ is discontinuous and a
pointwise comparison is ill-posed.}
\label{fig:twolayer}
\end{figure}

\section{Benchmark 2: DDH03 Figure 7 --- borehole, invasion, dipping
anisotropic layer}

\paragraph{Model (DDH03 Fig.~6).} The paper's flagship 3-D inhomogeneous
model, fully anisotropic and with no closed-form solution:
\begin{center}
\begin{tabular}{ll}
\toprule
Borehole & $\sigma = 0.05$ S/m, radius $0.1$ m \\
Invasion zone & $\sigma = 0.10$ S/m, radius $0.6$ m \\
Dipping anisotropic layer & $\sigma_T = 0.1$, $\sigma_N = 0.01$ S/m,
dip $60^\circ$, crossing the axis at $z=-0.5$ m \\
Isotropic layer & $\sigma = 0.5$ S/m \\
Source & $x$-directed magnetic dipole, $52.65$ kHz, at the origin \\
\bottomrule
\end{tabular}
\end{center}
The conductivity tensor of the dipping layer has its principal axes tilted
$60^\circ$ from vertical, so all off-diagonal couplings are active --- the
case Lebedev's grid exists for. Cells cut by the borehole wall, invasion
front, or dip plane are homogenized by sequential nodal homogenization with
analytical interface normals.

\paragraph{Test.} As in the paper, this is a \emph{self-consistency
convergence} test: the model is solved on optimal transverse grids with
$k = 3, 4, 5, 6$ steps per half-axis ($M_x = M_y = 4k$; hybrid axial grid
with $M_z = 98$; from $\approx31$k to $\approx123$k degrees of freedom), and
the on-axis magnetic induction components $\mathrm{Im}\,B_{xx}$ and
$\mathrm{Im}\,B_{xz}$ are compared across $k$. DDH03 report (p.~1535) that
their curves for $k = 4, 5, 6$ are ``practically identical'' with $k=3$
``almost as good.''

\paragraph{Results.} Figure~\ref{fig:fig7} shows the reproduction; the four
curves overlay to plotting accuracy for both components, including the
off-diagonal $B_{xz}$, which vanishes identically in a homogeneous medium and
is generated here purely by the dipping inhomogeneity. Quantitatively,
against the $k=6$ reference:

\begin{center}
\begin{tabular}{lcccc}
\toprule
 & \multicolumn{2}{c}{$\mathrm{Im}\,B_{xx}$}
 & \multicolumn{2}{c}{$\mathrm{Im}\,B_{xz}$} \\
 vs.\ $k=6$ & max & mean & max & mean \\
\midrule
$k=3$ \ ($M_x=12$) & 7.8\% & 4.4\% & 5.4\% & 3.3\% \\
$k=4$ \ ($M_x=16$) & 3.6\% & 2.2\% & 4.1\% & 3.1\% \\
$k=5$ \ ($M_x=20$) & 1.5\% & 0.8\% & 1.5\% & 1.2\% \\
\bottomrule
\end{tabular}
\end{center}

\noindent The deviations shrink monotonically with $k$, consistent with the
exponential convergence the optimal grids are designed to deliver, and the
$k=3$ solution --- only three geometric steps to each side of the borehole,
$M_x = 12$ --- is already within a few percent everywhere, mirroring the
paper's central claim that the scheme handles complicated 3-D anisotropic
models with remarkably few transverse nodes. Solve times on this hardware
ranged from 4 s ($k=3$) to 118 s ($k=6$) with a sparse direct solver.

\begin{figure}[htbp]
\centering
\includegraphics[width=\textwidth]{../examples/ddh03_fig7.png}
\caption{Reproduction of the DDH03 Figure~7 self-consistency test.
$\mathrm{Im}\,B_{xx}$ (left) and the dip-induced off-diagonal
$\mathrm{Im}\,B_{xz}$ (right) along the borehole axis for optimal grids with
$k=3,4,5,6$; the interface of the dipping anisotropic layer crosses the axis
at $z = -0.5$ m (dotted line).}
\label{fig:fig7}
\end{figure}

\section{Remarks}

Both benchmarks were run on the repository state including the July 2026
fixes (correct magnetic boundary rows in the curl operator, coordinate-based
cluster interpolation, corrected geometry straddle detection --- the latter
directly affects which near-axis cells are homogenized in Benchmark~2). The
two benchmarks are complementary: Benchmark~1 provides an \emph{external}
analytic reference but a simple geometry; Benchmark~2 has no analytic
solution but exercises the full anisotropic, multi-material machinery and
confirms internal convergence. A natural next step is the 2C-40 sonde
comparison of DDH03 Figure~5, for which quasi-analytic reference curves
exist in the literature (Anderson et al.~2001).

\begin{thebibliography}{9}
\bibitem{DDH03} S.~Davydycheva, V.~Druskin, T.~Habashy,
\emph{An efficient finite-difference scheme for electromagnetic logging in 3D
anisotropic inhomogeneous media}, Geophysics \textbf{68}(5):1525--1536, 2003.
\bibitem{MDHLD99} S.~Moskow, V.~Druskin, T.~Habashy, P.~Lee,
S.~Davydycheva, \emph{A finite difference scheme for elliptic equations with
rough coefficients using a Cartesian grid nonconforming to interfaces},
SIAM J.\ Numer.\ Anal.\ \textbf{36}:442--464, 1999.
\end{thebibliography}

\end{document}
